% modified from NeurIPS 2022 Latex template
% authors: Yixin Zhu (yixin.zhu@pku.edu.cn), Liangru Xiang
% v1.0: 2022.07.21

\documentclass{article}
\usepackage[final,nonatbib]{template22}
% Common packages
\usepackage{xeCJK}
\setmainfont{texgyretermes-regular.otf}[BoldFont=texgyretermes-bold.otf,ItalicFont=texgyretermes-italic.otf,BoldItalicFont=texgyretermes-bolditalic.otf]
\setsansfont{texgyreheros-regular.otf}[BoldFont=texgyreheros-bold.otf,ItalicFont=texgyreheros-italic.otf,BoldItalicFont=texgyreheros-bolditalic.otf]
\setmonofont{texgyrecursor-regular.otf}[BoldFont=texgyrecursor-bold.otf,ItalicFont=texgyrecursor-italic.otf,BoldItalicFont=texgyrecursor-bolditalic.otf]
\defaultCJKfontfeatures{}
\setCJKmainfont{Songti SC}
\setCJKsansfont{Songti SC}
\setCJKmonofont{Songti SC}
\usepackage[nopatch=footnote]{microtype}
\usepackage{graphicx}
\usepackage{amsmath,amssymb,mathbbol}
% \usepackage{algorithmic}
% \usepackage[linesnumbered,ruled,vlined]{algorithm2e}
\usepackage{acronym}
\usepackage{enumitem}
\usepackage[breaklinks=true,colorlinks]{hyperref}
\usepackage{balance}
\usepackage{xspace}
\usepackage{setspace}
\usepackage[skip=3pt,font=small]{subcaption}
\usepackage[skip=3pt,font=small]{caption}
\captionsetup[figure]{labelfont=bf,width=0.8\textwidth}
\captionsetup[subfigure]{width=\linewidth}
\usepackage[dvipsnames]{xcolor}
\usepackage[capitalise]{cleveref}
\usepackage{booktabs,tabularx,colortbl,multirow,array,makecell}
% \usepackage{overpic,wrapfig}

\usepackage{fancyhdr}
% Slightly widen the template's text area to keep the main text within four pages.
\AtBeginDocument{%
  \newgeometry{%
    left=1.25in,right=1.25in,top=1in,textheight=9in,
    headheight=12pt,headsep=25pt,footskip=30pt
  }%
}
\hypersetup{pdfencoding=auto,colorlinks=true,allcolors=black}
\renewcommand{\headrulewidth}{0.5pt}
\renewcommand{\footrulewidth}{0pt}
\fancyhf{}
\fancyhead[C]{}
\fancyhead[C]{}
\fancyfoot[C]{\thepage}

% Handy shorthand
\makeatletter
\DeclareRobustCommand\onedot{\futurelet\@let@token\@onedot}
\def\@onedot{\ifx\@let@token.\else.\null\fi\xspace}
\def\eg{\emph{e.g}\onedot}
\def\Eg{\emph{E.g}\onedot}
\def\ie{\emph{i.e}\onedot}
\def\Ie{\emph{I.e}\onedot}
\def\cf{\emph{c.f}\onedot}
\def\Cf{\emph{C.f}\onedot}
\def\etc{\emph{etc}\onedot}
\def\vs{\emph{vs}\onedot}
\def\wrt{w.r.t\onedot}
\def\dof{d.o.f\onedot}
\def\etal{\emph{et al}\onedot}
\makeatother

\definecolor{gray}{gray}{0.9}

% Handy math ops
\DeclareMathOperator*{\argmax}{arg\,max}
\DeclareMathOperator*{\argmin}{arg\,min}
\newcommand{\norm}[1]{\left\Vert #1 \right\Vert}

% % Spacing
\frenchspacing
% \medmuskip=2mu   % reduce spacing around binary operators
% \thickmuskip=3mu % reduce spacing around relational operators
% \setlength{\abovedisplayskip}{3pt}
% \setlength{\belowdisplayskip}{3pt}
% \setlength{\abovecaptionskip}{3pt}
% \setlength{\belowcaptionskip}{3pt}
\setlength\floatsep{0.5\baselineskip plus 3pt minus 2pt}
\setlength\textfloatsep{0.5\baselineskip plus 3pt minus 2pt}
\setlength\dbltextfloatsep{0.5\baselineskip plus 3pt minus 2pt}
\setlength\intextsep{0.5\baselineskip plus 3pt minus 2pt}

\makeatletter
\renewcommand{\paragraph}{%
  \@startsection{paragraph}{4}%
  {\z@}{0ex \@plus 0ex \@minus 0ex}{-1em}%
  {\hskip\parindent\normalfont\normalsize\bfseries}%
}
\makeatother

% Graphics path
\graphicspath{{figures/}}

% Clever references
\crefname{algorithm}{Alg.}{Algs.}
\Crefname{algorithm}{Algorithm}{Algorithms}
\crefname{section}{Sec.}{Secs.}
\Crefname{section}{Section}{Sections}
\crefname{table}{Tab.}{Tabs.}
\Crefname{table}{Table}{Tables}
\crefname{figure}{Fig.}{Fig.}
\Crefname{figure}{Figure}{Figure}

% Acronym
\acrodef{pku}[PKU]{Peking University}

\usepackage{listings}
\definecolor{pythonkeyword}{HTML}{7C3AED}
\definecolor{pythonstring}{HTML}{047857}
\definecolor{pythoncomment}{HTML}{64748B}
\definecolor{pythonbackground}{HTML}{FAF6ED}
\lstdefinestyle{python}{
  language=Python,
  basicstyle=\small\ttfamily,
  backgroundcolor=\color{pythonbackground},
  keywordstyle=\color{pythonkeyword}\bfseries,
  stringstyle=\color{pythonstring},
  commentstyle=\color{pythoncomment}\itshape,
  morekeywords={print,chr},
  showstringspaces=false,
  columns=fullflexible,
  breaklines=true
}
\usepackage[style=apa,backend=biber,natbib=true,backref=false]{biblatex}
\addbibresource{reference.bib}
\newif\ifhascitations
\AtEveryCitekey{\global\hascitationstrue}

\title{Homework 1: Byte-Pair Encoding and Language Model}

\author{%
  Koh Zhi Xiang\\
  Yuanpei College\\
  Peking University\\
  \texttt{2300094603@stu.pku.edu.cn} \\
}

\begin{document}
\maketitle

\section*{2 Byte-Pair Encoding (BPE) Tokenizer}
\label{sec:bpe}

\subsection*{2.1 Understanding Unicode}

\begin{enumerate}[label=(\alph*), leftmargin=*]
  \item \texttt{\textbackslash x00} (NUL byte).
  \item The printed representation is empty compared to that of subproblem (a).
  \item
\begin{lstlisting}[style=python]
>>> "this is a test" + chr(0) + "string"
'this is a test\x00string'
>>> print("this is a test" + chr(0) + "string")
this is a teststring
\end{lstlisting}
\end{enumerate}

\subsection*{2.2 Unicode Encodings}

\begin{enumerate}[label=(\alph*), leftmargin=*]
  \item UTF-16/32 introduce redundant \texttt{\textbackslash x00} byte when encoding common characters (prefix in BE variant), resulting in large byte sequences than UTF-8. This increase storage and preprocessing costs.
  \item The function did not consider characters encoded by more than one byte (\texttt{chr(x)} where \texttt{x} is larger than 127). It returns \textit{'utf-8' codec can't decode byte 0xc2 in position 0: unexpected end of data} \texttt{UnicodeDecodeError}.
  \item \texttt{\textbackslash x00 \textbackslash x00}
\end{enumerate}

\subsection*{2.3 Subword Tokenization and BPE Tokeniser}

\subsection*{2.4 Experimenting with BPE Tokenizer Training}

\section*{4 LM Training}

\subsection*{4.2 Implementing SGD in PyTorch}

\begin{table}[htbp]
  \centering
  \small
  \caption{Training loss at each iteration for three initial learning rates.}
  \label{tab:sgd-learning-rates}
  \begin{tabular}{@{}rrrr@{}}
    \toprule
    Iteration $t$ & $\alpha = 10$ & $\alpha = 100$ & $\alpha = 1000$ \\
    \midrule
    0 & 21.39 & 20.53 & 26.84 \\
    1 & 13.69 & 20.53 & 9689 \\
    2 & 10.09 & 3.523 & $1.673 \times 10^{6}$ \\
    3 & 7.895 & 0.08432 & $1.861 \times 10^{8}$ \\
    4 & 6.395 & $1.595 \times 10^{-16}$ & $1.508 \times 10^{10}$ \\
    5 & 5.302 & $1.778 \times 10^{-18}$ & $9.516 \times 10^{11}$ \\
    6 & 4.472 & $5.988 \times 10^{-20}$ & $4.885 \times 10^{13}$ \\
    7 & 3.821 & $3.567 \times 10^{-21}$ & $2.102 \times 10^{15}$ \\
    8 & 3.300 & $3.060 \times 10^{-22}$ & $7.747 \times 10^{16}$ \\
    9 & 2.875 & $3.400 \times 10^{-23}$ & $2.488 \times 10^{18}$ \\
    \bottomrule
  \end{tabular}
\end{table}

Clearly, this shows that $\alpha=100$ converges faster than $\alpha=10$, while $\alpha=1000$ diverges.

\clearpage
\subsection*{4.3 AdamW}

Let $B,T,V,L,d,h$ denote batch size, context length, vocabulary size, layer count, model width and head count. Set $d_{\mathrm{ff}}=4d$. Each float32 element occupies 4 bytes; $1\,\mathrm{GB}=10^9$ bytes.

\paragraph{(a) Memory.}
The two embedding matrices, attention projections, three SwiGLU matrices and RMSNorm gains give
\[
P=2Vd+L\bigl(4d^2+12d^2+2d\bigr)+d=2Vd+L(16d^2+2d)+d.
\]
Parameter, gradient and Adam-state memory are, respectively,
\[
M_{\mathrm{param}}=4P,\qquad M_{\mathrm{grad}}=4P,\qquad M_{\mathrm{Adam}}=8P
\quad\text{bytes}.
\]
Counting the listed activation outputs once gives:
\begin{center}
\small
\begin{tabular}{@{}lr@{}}
\toprule
Component & Float32 elements \\
\midrule
Two RMSNorms per block & $2BTd$ \\
$\mathbf{Q}$, $\mathbf{K}$, $\mathbf{V}$, weighted values and output projection per block & $5BTd$ \\
Attention scores and softmax probabilities per block & $2BhT^2$ \\
SwiGLU: $\mathbf{W}_1\mathbf{x}$, SiLU, $\mathbf{W}_3\mathbf{x}$, gating product and $\mathbf{W}_2$ output & $17BTd$ \\
\midrule
Final RMSNorm & $BTd$ \\
Logits & $BTV$ \\
Cross-entropy: negative log probabilities and token losses & $BTV+BT$ \\
\bottomrule
\end{tabular}
\end{center}
Thus $A=BT[L(24d+2hT)+d+2V+1]$ and $M_{\mathrm{act}}=4A$ bytes, giving
\[
M_{\mathrm{peak}}\approx16P+4A\quad\text{bytes}.
\]

\paragraph{(b) GPT-2 XL.}
Using $(V,T,L,d,h)=(50{,}257,1024,48,1600,25)$,\footnote{Model configuration: \url{https://huggingface.co/openai-community/gpt2-xl/blob/main/config.json}.}
\[
P=2{,}127{,}057{,}600,\qquad A=4{,}508{,}584{,}960\,B.
\]
Parameters and gradients each occupy $8.508$ GB, and Adam states occupy $17.02$ GB. The total is
\begin{align*}
M_{\mathrm{peak}}(B)
 &= (18.03433984B+34.0329216)\,\mathrm{GB}\approx(18.03B+34.03)\,\mathrm{GB},\\
B_{\max}
 &=\left\lfloor\frac{80-34.0329216}{18.03433984}\right\rfloor
 =2.
\end{align*}
Batch sizes 2 and 3 require $70.10$ GB and $88.14$ GB, respectively.

\paragraph{(c) AdamW FLOPs.}
Count each scalar arithmetic operation and square root as one FLOP, with scalar coefficients precomputed. Per parameter, updating $\mathbf{m}$ costs 3 FLOPs, updating $\mathbf{v}$ costs 4, the adaptive parameter update costs 5, and weight decay costs 2. Hence
\[
F_{\mathrm{AdamW}}\approx(3+4+5+2)P=14P
\]
FLOPs per update, giving $2.978\times10^{10}$ FLOPs for (b).

\paragraph{(d) Training time.}
Count dense matrix products at $2mkn$ FLOPs. Per block, $\mathbf{Q}$, $\mathbf{K}$, $\mathbf{V}$ and output projections cost $8BTd^2$ FLOPs, attention products $4BT^2d$, the three SwiGLU projections $24BTd^2$, and gating $4BTd$. Adding the vocabulary projection gives
\[
F_{\mathrm{fwd}}\approx BT\bigl[L(32d^2+4Td+4d)+2dV\bigr].
\]
For $B=1024$, $F_{\mathrm{fwd}}\approx4.622\times10^{15}$ FLOPs. Backward costs twice forward, so each training step requires approximately $3F_{\mathrm{fwd}}=1.387\times10^{16}$ model FLOPs. The effective rate at 50\% MFU is $9.75\times10^{12}$ FLOPs/s. Thus
\[
\text{time}\approx
\frac{400{,}000\times3F_{\mathrm{fwd}}}{9.75\times10^{12}\times86{,}400}
=6.584\times10^3\ \text{days}\approx18.04\ \text{years}.
\]
An effective batch of 1024 can use gradient accumulation over 512 microbatches of size 2.

\ifhascitations\printbibliography\fi

\end{document}
